%% encoding: UTF-8
%% 首先读取模板的配置
%%%%%%%%%%%%%%%%%%%%%%%%%%%%%%%设置版面控制%%%%%%%%%%%%%%%%%%%%%%%%%%%%%%%%%%
\documentclass[a4paper, 10pt, notitlepage, oneside , twoside ]{article}
\usepackage[colorlinks, linkcolor=red, anchorcolor=blue, citecolor=green, bookmarks]{hyperref}
\hypersetup{xetex}
%\setlength{\parindent}{0em}
\usepackage[papersize={158.8mm,201.3mm},top=0cm,bottom=0.5cm,left=0cm,right=0cm]{geometry}
%\usepackage[papersize={200mm,255mm},top=3.37cm,bottom=2cm,left=2.06cm,right=2.06cm]{geometry}
%\usepackage[papersize={210mm,297mm},top=2.8cm,bottom=2.5cm,left=2.5cm,right=2.5cm]{geometry}                 % 设置版面，纸张空距
\usepackage{marvosym}
\usepackage{makecell,multirow}

%%%%%%%%%%%%%%%%%%%%%%%%%%%%%设置字体%%%%%%%%%%%%%%%%%%%%%%%%%%%%%
\usepackage[no-math]{fontspec} %提供字体选择命令
\usepackage{xunicode,xltxtra}
\usepackage[CJKchecksingle, CJKnumber]{xeCJK}                 % 用于字体
%\setCJKmainfont[BoldFont={FZLanTingHei-DB-GBK}]{Microsoft YaHei}
\setCJKmainfont[BoldFont={FZLanTingHei-DB-GBK}]{FZLanTingHei-L-GBK}
%\setmainfont{Microsoft YaHei}   % 设置中文/粗体/斜体/主要字体
\setmainfont{FZLanTingHei-L-GBK}
\punctstyle{hangmobanjiao}               % 设置行末半角
\newfontfamily\times{Times New Roman}
\newCJKfontfamily\kaishu{Adobe Kaiti Std}
\setCJKfamilyfont{zhyh}{Microsoft YaHei}
\newcommand*{\heiti}{\CJKfamily{zhyh}}
\setCJKfamilyfont{zhltxh}{FZLanTingHei-L-GBK}
\newcommand*{\lantxh}{\CJKfamily{zhltxh}}
\newfontfamily\arial{Arial}
\newcommand{\barial}{\bf\arial}
%\setmainfont[Mapping=tex-text]{Times New Roman}         % 设置衬线英文字体
\setsansfont[Mapping=tex-text]{FZLanTingHei-L-GBK}       % 设置无衬线英文字体
%\setmonofont{Courier New}                               % 设置等款英文字体
\newcommand{\chuhao}{\fontsize{42pt}{\baselineskip}\selectfont}
\newcommand{\xiaochuhao}{\fontsize{36pt}{\baselineskip}\selectfont}
\newcommand{\yihao}{\fontsize{28pt}{\baselineskip}\selectfont}
\newcommand{\erhao}{\fontsize{21pt}{\baselineskip}\selectfont}
\newcommand{\xiaoerhao}{\fontsize{18pt}{\baselineskip}\selectfont}
\newcommand{\sanhao}{\fontsize{15.75pt}{\baselineskip}\selectfont}
\newcommand{\sihao}{\fontsize{14pt}{\baselineskip}\selectfont}
\newcommand{\xiaosihao}{\fontsize{12pt}{\baselineskip}\selectfont}
\newcommand{\wuhao}{\fontsize{10.5pt}{\baselineskip}\selectfont}
\newcommand{\xiaowuhao}{\fontsize{9pt}{15pt}\selectfont}
\newcommand{\liuhao}{\fontsize{7.875pt}{\baselineskip}\selectfont}
\newcommand{\qihao}{\fontsize{7pt}{\baselineskip}\selectfont}
\newcommand{\bahao}{\fontsize{6.875pt}{\baselineskip}\selectfont}
\newcommand{\wudianwu}{\fontsize{5.5pt}{\baselineskip}\selectfont}
\newcommand{\RNum}[1]{\uppercase\expandafter{\romannumeral #1\relax}}
\newcommand{\rnum}[1]{\lowercase\expandafter{\romannumeral #1\relax}}

\usepackage{textgreek}
\usepackage{pdflscape}
%%%%%%%%%%%%%%%%%%%%%%%%%%%%%设置图像%%%%%%%%%%%%%%%%%%%%%%%%%%%%%
\usepackage[usenames,dvipsnames,svgnames,table]{xcolor}
\usepackage{graphicx}                    % 用于图像编号
\graphicspath{{cores/}}
\usepackage{tikz}
\usetikzlibrary{calc}
\usetikzlibrary{shapes.geometric, arrows}
\usetikzlibrary{circuits.logic.US,circuits.logic.IEC,circuits.ee.IEC}
\usepackage{verbatim}
\usepackage{animate}
\usepackage{subfig}
%\usepackage{mathpazo}
\usepackage{wrapfig}
\usepackage{pgfplots}
\pgfplotsset{compat=newest}% <-- moves axis labels near ticklabels (respects tick label widths)
\usepackage{wallpaper}%%%%%%%%%%%%%%%%背景%%%%%%%%%%%%%%%%%%
\usepackage{color}
%%%%%%%%%%%%%%%%%%%%%%%%%%mdframed%%%%%%%%%%%%%%%%%%%%
\usepackage[framemethod=tikz]{mdframed}
\usetikzlibrary{patterns}
\mdfdefinestyle{titlesty}{%
linecolor=platinum,
middlelinewidth=1.5pt,
frametitlerule=false,%
apptotikzsetting={\tikzset{
    mdfbackground/.append style ={
    top color=white,
    bottom color=white
    },
    mdfframetitlebackground/.append style={
    shade,
    top color=white ,
    bottom color=white
    }
}
},
frametitlerulewidth=0pt,
innertopmargin=\topskip,
topline=false,
leftline=false,
rightline=false,
bottomline=false,
frametitlealignment=\raggedright,
}
\usepackage{fancyhdr}
\pagestyle{fancy}
\fancyhf{}
\renewcommand{\headrulewidth}{0pt}
\renewcommand{\footrulewidth}{0pt}
%%%%%%%%%%%%%%%页码修改[RO]居右[CO]居中[LO]居左%%%%%%%%%%%

%\cfoot{}
%\fancyfoot[LO,RE]{}
%\fancyhead[LO,RE]{\raisebox{-1.5ex}{\includegraphics{logo}}\color{gray} ~\bigg|\,~\makecell[l]{胃动力~~肠健康\\
%\scalebox{0.8}{Better Gut Better Life!}}}

\usepackage{color}
\usepackage{graphicx}
\usepackage{paralist}
\usepackage{float,longtable,hhline}

\definecolor{topcolor}{cmyk}{0.28,0.00,1.00,0.00}

\definecolor{solomon0}{cmyk}{0.77, 0.53, 0.00, 0.00}
\definecolor{solomon1}{cmyk}{0.62, 0.00, 0.12, 0.00}
\definecolor{solomon2}{cmyk}{0.43, 0.67, 0.00, 0.36}
\definecolor{solomon3}{cmyk}{1.00, 0.18, 0.40, 0.16}
\definecolor{solomon4}{cmyk}{0.26, 0.00, 0.00, 0.00}
\definecolor{solomon5}{cmyk}{0.80, 0.30, 0.05, 0.32}
\definecolor{solomon6}{cmyk}{0.32, 0.24, 0.00, 0.00}
\definecolor{xxx1}{RGB}{46,157,148}
\definecolor{green}{cmyk}{0.3,0.03,0.39,0}
%\definecolor{green2}{RGB}{196,224,177}
\definecolor{green2}{cmyk}{0.3,0.03,0.39,0}
\definecolor{yellow}{RGB}{238,196,133}
\definecolor{yellow1}{RGB}{243,225,0}
\definecolor{columncolor}{cmyk}{0,0,0,.08}%{RGB}{247,248,248}
\definecolor{red}{RGB}{225,135,117}
\definecolor{orange}{RGB}{245,130,31}
\definecolor{gray1}{RGB}{210,210,210}
%\definecolor{gray2}{RGB}{144,144,144}
\definecolor{gray0}{RGB}{244,244,244}
\definecolor{blue1}{RGB}{65,100,140}
\definecolor{blue2}{RGB}{50,168,225}
\definecolor{kelly}{cmyk}{0.28,0.00,1.00,0.00}
\definecolor{black70}{cmyk}{0.00,0.00,0.00,0.90}
\definecolor{darkblue}{cmyk}{0.74,0.02,0.29,0.00}
\definecolor{gray2}{cmyk}{0.00,0.00,0.00,0.90}

\def\blueC{{\color{blue}C}}
\def\blueG{{\color{blue}G}}
\def\redA{{\color{red}A}}
\def\redT{{\color{red}T}}
\def\yellowA{{\color{yellow}A}}
\def\yellowT{{\color{yellow}T}}
\def\redrule{\raisebox{.5ex}{~~~\hspace*{-3cm}{\color{red}\rule{2.3cm}{1.5pt}}}~}

\def\sun{\tikz\draw[yellow1,fill=yellow1](0,0) circle[radius=4.5pt];
\hspace*{-0.5cm}\tikz\draw[yellow1,dotted](0,0) circle[radius=5pt];}


\def\trianglea{\tikz \draw[scale=0.07,color=yellow1,fill=yellow1] (1,0) -- (-1,0) -- (0,1.732) -- cycle;}
\def\triangleb{\tikz \draw[scale=0.07,color=yellow1,fill=yellow1] (1,0) -- (-1,0) -- (0,-1.732) -- cycle;}
\usepackage{tabularx}

\usepackage[many]{tcolorbox}
%\newtcolorbox{tctabularx}[1]{fonttitle=\bfseries, colback=white, colframe=kelly,#1,sharp corners,rounded corners=northwest,before upper app={\rowcolor{kelly}}}   % 表头和minipage左上角设置为圆角
\newtcolorbox{tctabularx}[1]{fonttitle=\bfseries, colback=white, colframe=kelly,#1,sharp corners,before upper app={\rowcolor{kelly}}}    % 表头和minipage左上角设置为直角

\makeatletter
\tcbset
    {
tabularx*/.style 2 args={%
    boxsep=0pt,top=0pt,bottom=0pt,leftupper=0pt,rightupper=0pt,
    toptitle=1mm,bottomtitle=1mm,boxrule=0.5mm,
    before upper={\arrayrulecolor{black70}\def\arraystretch{1}#1%
      \tcb@hack@currenvir\tabularx{\linewidth}{#2}},
    after upper=\endtabularx\arrayrulecolor{gray}}
    }
\makeatother
